\documentclass[11pt]{article}

\usepackage[margin=1in]{geometry}
\usepackage[T1]{fontenc}
\usepackage[utf8]{inputenc}
\usepackage{lmodern}
\usepackage{amsmath,amssymb,mathtools}
\usepackage{booktabs}
\usepackage{longtable}
\usepackage{array}
\usepackage{enumitem}
\usepackage{xcolor}
\usepackage[numbers,sort&compress]{natbib}
\usepackage{hyperref}
\usepackage{url}

\hypersetup{
  colorlinks=true,
  linkcolor=blue!50!black,
  citecolor=blue!50!black,
  urlcolor=blue!50!black
}

\newcommand{\R}{\mathbb{R}}
\newcommand{\DeltaN}{\Delta}
\newcommand{\ones}{\mathbf{1}}
\newcommand{\diag}{\operatorname{diag}}
\newcommand{\rankp}{\operatorname{rank}_{+}}
\newcommand{\KL}{\operatorname{KL}}
\newcommand{\OT}{\operatorname{OT}}
\newcommand{\SW}{\operatorname{SW}}
\newcommand{\MSW}{\operatorname{MSW}}
\newcommand{\push}{_{\#}}
\newcommand{\inner}[2]{\left\langle #1,#2\right\rangle}

\title{Scalable Optimal Transport for LEDH-PFPF-OT:\\
A Self-Contained Literature Survey}
\author{BayesFilter local literature survey}
\date{June 17, 2026}

\begin{document}
\maketitle

\begin{abstract}
LEDH-PFPF-OT combines an invertible particle-flow proposal with an optimal
transport or differentiable-resampling step.  The scaling problem is not merely
to compute an optimal-transport distance.  The filter needs a transport object:
a coupling, a barycentric projection, or a directly transported particle cloud
that maps a weighted empirical ensemble to an approximately equal-weight
ensemble.  Dense entropic optimal transport provides such an object, but for
$N$ particles in state dimension $D$ it requires all-pairs costs or kernels:
$O(N^2D)$ arithmetic per full squared-Euclidean pass and $O(N^2)$ storage if
the cost, kernel, or coupling is materialized.  The current TensorFlow
BayesFilter implementation already has a streaming mode that avoids storing
the full transport matrix, but it still computes all row--column particle
interactions.

This survey explains the LEDH-PFPF-OT transport problem and reviews the main
solution families in the scalable optimal-transport literature: exact online
and GPU Sinkhorn methods, Nystrom kernel Sinkhorn, positive-feature Sinkhorn,
direct low-rank coupling optimal transport, sparse and screened methods,
accelerated Sinkhorn variants, stochastic and minibatch optimal transport,
sliced and subspace optimal transport, and particle-filter-specific
localization.  The recommended research posture is a three-lane program:
retain exact online Sinkhorn as the semantics-preserving reference, prototype
factored entropic transports such as Nystrom and positive-feature Sinkhorn,
and study low-rank coupling, localized/block, and sliced/subspace transports
as explicit new resampling methods requiring downstream filtering validation.
\end{abstract}

\section{Evidence Contract and Scope}

\begin{center}
\begin{tabular}{p{0.20\textwidth}p{0.72\textwidth}}
\toprule
Field & Contract \\
\midrule
Question & How can the dense or all-pairs optimal-transport step in
LEDH-PFPF-OT be scaled for large particle count $N$ and large state dimension
$D$? \\
Baseline & Current BayesFilter
\path{annealed_transport_resample_tf} in
\path{experiments/dpf_implementation/tf_tfp/resampling/annealed_transport_tf.py}. \\
Required object & A usable particle transport: a coupling, a barycentric
projection, or a directly transported cloud. \\
Screen & Does the method reduce $N\times N$ memory or all-pairs
$O(N^2D)$ arithmetic while preserving or explicitly replacing the transport
object? \\
Veto & A method that returns only a scalar loss; cannot handle weighted
empirical measures; silently changes resampling semantics; or depends on a
non-TensorFlow backend as the BayesFilter implementation default. \\
Not concluded & No method is default-ready from literature alone.  No
posterior-validity, gradient-correctness, HMC-readiness, or statistically
supported ranking claim is made in this document. \\
\bottomrule
\end{tabular}
\end{center}

The ResearchAssistant workspace at \path{/home/ubuntu/python/ResearchAssistant}
was available through a read-only MCP adapter, but its local index returned no
stored summaries for this topic.  The evidence base is therefore the local PDF,
text, and source corpus listed in Appendix~\ref{sec:source-map}.  A later arXiv
download attempt for one additional 2025 source stalled and produced an
incomplete file; that partial file is recorded in the manifest and is not used
as evidence here.

\section{The LEDH-PFPF-OT Computation Problem}

\subsection{Particle-flow context}

The Li--Coates PF-PF(LEDH) construction used in the local BayesFilter
documentation is a proposal-corrected particle-flow particle filter
\citep{li2017particle}; the broader particle-flow filtering idea follows the
Daum--Huang particle-flow line \citep{daumhuang2008}.  It is not
merely a deterministic cloud motion.  For particle $i$, the LEDH flow maintains
both a zero-noise auxiliary path and an actual proposal path.  With pseudo-time
$\lambda$, the local observation linearization is evaluated at the moving
auxiliary state:
\begin{equation}
H^i(\lambda)
=
\left.
\frac{\partial h(\eta,0)}{\partial \eta}
\right|_{\eta=\bar\eta^i_\lambda},
\qquad
e^i(\lambda)
=
h(\bar\eta^i_\lambda,0)-H^i(\lambda)\bar\eta^i_\lambda .
\label{eq:ledh-local-jacobian}
\end{equation}
Given particle-specific predicted covariance $P^i$, the local LEDH coefficients
are
\begin{align}
A^i(\lambda)
&=
-\frac{1}{2}
P^i H^i(\lambda)^\top
\bigl(\lambda H^i(\lambda)P^iH^i(\lambda)^\top+R\bigr)^{-1}
H^i(\lambda),
\label{eq:ledh-A}\\
b^i(\lambda)
&=
\bigl(I+2\lambda A^i(\lambda)\bigr)
\left[
  \bigl(I+\lambda A^i(\lambda)\bigr)
  P^i H^i(\lambda)^\top R^{-1}\{z_k-e^i(\lambda)\}
  + A^i(\lambda)\bar\eta_0^i
\right].
\label{eq:ledh-b}
\end{align}
On a pseudo-time grid
$0=\lambda_0<\lambda_1<\cdots<\lambda_{N_\lambda}=1$, the auxiliary and actual
states are updated by the same local affine step:
\begin{align}
\bar\eta^i_{\lambda_j}
&=
\bar\eta^i_{\lambda_{j-1}}
+ \epsilon_j
\left[
  A_j^i(\lambda_j)\bar\eta^i_{\lambda_{j-1}}
  + b_j^i(\lambda_j)
\right],
\label{eq:aux-step}\\
\eta^i_{\lambda_j}
&=
\eta^i_{\lambda_{j-1}}
+ \epsilon_j
\left[
  A_j^i(\lambda_j)\eta^i_{\lambda_{j-1}}
  + b_j^i(\lambda_j)
\right].
\label{eq:actual-step}
\end{align}
The forward determinant multiplier is particle-specific:
\begin{equation}
\theta^i
=
\prod_{j=1}^{N_\lambda}
\left|\det\{I+\epsilon_j A_j^i(\lambda_j)\}\right|.
\label{eq:theta-product}
\end{equation}
After the flow, Algorithm 1 uses a PF-PF importance update of the form
\begin{equation}
w_k^i
\propto
\frac{
  p(x_k^i\mid x_{k-1}^i)
  p(z_k\mid x_k^i)
  \theta^i
}{
  p(\eta_0^i\mid x_{k-1}^i)
}
w_{k-1}^i .
\label{eq:pfpf-weight}
\end{equation}
If resampling is used, the documented Algorithm 1 state to resample is the
triple
\begin{equation}
\{x_k^i,P_k^i,w_k^i\}_{i=1}^{N_p}.
\label{eq:resampling-triple}
\end{equation}
Thus optimal-transport resampling is a BayesFilter extension layered on top of
a proposal-corrected flow; the OT-resampling motivation is closer to ensemble
transform and differentiable OT-resampling work
\citep{reich2013nonparametric,corenflos2021differentiable}.  The survey
question is how to make that extension scalable without confusing it with the
Li--Coates source algorithm \citep{li2017particle}.

\subsection{Why a scalar OT value is not enough}

At a filtering time step, suppose the flow and likelihood update have produced
particles $x_i\in\R^D$ with normalized weights $a_i$.  A deterministic
OT-resampling step compares the weighted empirical measure, in the spirit of
ensemble-transform and differentiable OT resampling
\citep{reich2013nonparametric,corenflos2021differentiable},
\begin{equation}
\mu_N = \sum_{i=1}^N a_i \delta_{x_i}
\label{eq:source-empirical}
\end{equation}
to an equal-weight target ensemble.  In the simplest self-transport version,
the support points are the same cloud and the target weights are
\begin{equation}
b_j = \frac{1}{N},
\qquad
\nu_N = \sum_{j=1}^N b_j \delta_{x_j}.
\label{eq:target-empirical}
\end{equation}
A coupling $P$ assigns source mass from particle $i$ to target slot $j$.  The
transport polytope is
\begin{equation}
\Pi(a,b)
=
\{P\in\R_+^{N\times N}: P\ones=a,\;P^\top\ones=b\}.
\label{eq:transport-polytope}
\end{equation}
The deterministic barycentric transform is
\begin{equation}
z_j
=
\frac{1}{b_j}\sum_{i=1}^N P_{ij}x_i.
\label{eq:barycentric-column}
\end{equation}
If the implementation stores a row-stochastic target-to-source matrix $T$, the
same transform is
\begin{equation}
z_i = \sum_{j=1}^N T_{ij}x_j,
\qquad
\sum_j T_{ij}=1.
\label{eq:barycentric-row}
\end{equation}

This is the central computational requirement.  A fast scalar
$\OT(\mu,\nu)$ is insufficient for LEDH-PFPF-OT.  The filter needs $P$, a
factorization of $P$, or an operator that can apply $P$ to the particle matrix
$X$.  This agrees with the ensemble transform particle filter, where random
resampling is replaced by the deterministic linear transform
\citep{reich2013nonparametric}
\begin{equation}
x_j^a = \bar x_j^a = \sum_{i=1}^M p_{ij} x_i^f,
\qquad j=1,\ldots,M.
\label{eq:reich-transform}
\end{equation}
Differentiable OT resampling likewise recovers a regularized transport matrix
from dual potentials \citep{corenflos2021differentiable}:
\begin{equation}
p^{OT}_{\epsilon,i,j}
=
a_i b_j
\exp\{\epsilon^{-1}(f_i^\ast+g_j^\ast-c_{ij})\}.
\label{eq:corenflos-plan}
\end{equation}

\subsection{The current BayesFilter bottleneck}

The current TensorFlow implementation has dense and streaming transport modes.
Dense mode materializes a transport matrix and applies $Z=TX$.  Streaming mode
returns an empty transport-matrix sentinel but still loops over row and column
chunks.  For each block it computes squared costs such as
\begin{equation}
C_{ij} = \frac{1}{2}\|\tilde x_i-\tilde x_j\|_2^2
\label{eq:current-cost}
\end{equation}
and block log-transport weights of the schematic form
\begin{equation}
\log T_{ij}
=
\frac{f_i+g_j-C_{ij}}{\epsilon}
- \ell_j
+ \log N
+ \log w_j,
\label{eq:current-log-transport}
\end{equation}
where $\ell_j$ is a column log normalizer.  It then accumulates
\begin{equation}
z_i = \sum_j T_{ij}x_j.
\label{eq:current-streamed-application}
\end{equation}
Streaming therefore solves part of the memory problem, but not the all-pairs
arithmetic problem.  Each full Sinkhorn update or transport application still
depends on all pairs of particles.

\section{Discrete OT and Entropic Sinkhorn Basics}

This section fixes notation for standard discrete OT and entropic Sinkhorn
formulations \citep{villani2003topics,cuturi2013sinkhorn,peyre2019computational}.

Let
\begin{equation}
\mu=\sum_{i=1}^n a_i\delta_{x_i},
\qquad
\nu=\sum_{j=1}^m b_j\delta_{y_j},
\label{eq:general-empirical}
\end{equation}
with $a\in\Delta_n$, $b\in\Delta_m$, and cost matrix
\begin{equation}
C_{ij}=c(x_i,y_j).
\label{eq:cost-matrix}
\end{equation}
The Kantorovich problem is
\begin{equation}
\OT(\mu,\nu)
=
\min_{P\in\Pi(a,b)} \inner{C}{P},
\label{eq:kantorovich}
\end{equation}
where
\begin{equation}
\Pi(a,b)
=
\{P\in\R_+^{n\times m}:P\ones_m=a,\;P^\top\ones_n=b\}.
\label{eq:general-polytope}
\end{equation}
Entropic regularization replaces the linear program by
\citep{cuturi2013sinkhorn,peyre2019computational}
\begin{equation}
\OT_\epsilon(\mu,\nu)
=
\min_{P\in\Pi(a,b)}
\inner{C}{P}
-\epsilon H(P),
\label{eq:entropic-ot}
\end{equation}
with entropy convention
\begin{equation}
H(P)=-\sum_{ij}P_{ij}(\log P_{ij}-1).
\label{eq:entropy}
\end{equation}
A relative-entropy version often used in differentiable resampling is
\begin{equation}
W^2_{2,\epsilon}(\alpha_N,\beta_N)
=
\min_{P\in S(a,b)}
\sum_{i,j=1}^N
\left[
  p_{ij}c_{ij}
  + \epsilon p_{ij}\log\frac{p_{ij}}{a_i b_j}
\right].
\label{eq:relative-entropic-ot}
\end{equation}
For $\epsilon>0$, the minimizer is unique and has Gibbs scaling form
\begin{equation}
P_\epsilon
=
\diag(u)K\diag(v),
\qquad
K=\exp(-C/\epsilon).
\label{eq:gibbs-scaling}
\end{equation}
The basic Sinkhorn updates are
\begin{equation}
u \leftarrow a/(Kv),
\qquad
v \leftarrow b/(K^\top u),
\label{eq:sinkhorn-updates}
\end{equation}
with division elementwise.  Stabilized log-domain variants use dual potentials
and log-sum-exp reductions rather than explicitly forming kernel entries that
can underflow \citep{schmitzer2019stabilized}.

For particle transport, the coupling is applied as
\begin{equation}
Z
=
\diag(b)^{-1}P^\top X,
\label{eq:coupling-application}
\end{equation}
or, under the row-stochastic orientation, as $Z=TX$.  Every scalable method
must explain how it computes this transported particle matrix.

\section{Exact Online and GPU Sinkhorn}

Online Sinkhorn preserves the exact entropic OT problem but avoids storing the
full cost or kernel matrix, as in GeomLoss/KeOps-style reductions and recent
GPU-specialized Sinkhorn solvers
\citep{feydy2019geomloss,charlier2021keops,ye2026flashsinkhorn,xiao2026fastsinkhorn}.
For squared Euclidean costs,
\begin{equation}
C_{ij}=\|x_i-y_j\|_2^2
=\|x_i\|_2^2+\|y_j\|_2^2-2x_i^\top y_j.
\label{eq:sqeuclidean-dot}
\end{equation}
The stabilized update for one potential is a row-wise log-sum-exp reduction,
schematically
\begin{equation}
f_i
\leftarrow
-\epsilon\log\sum_j
\exp\left((g_j-C_{ij})/\epsilon+\log b_j\right),
\label{eq:log-domain-update}
\end{equation}
and similarly for $g_j$.  Online backends compute these reductions by
streaming blocks of $x_i,y_j$ rather than reading a stored $C$ or $K$.

FlashSinkhorn observes that, for squared Euclidean costs, the log-domain update
can be written as a biased dot-product log-sum-exp, close to transformer
attention normalization \citep{ye2026flashsinkhorn}.  It then uses tiled GPU
kernels to stream data through on-chip memory.  The paper is especially
relevant because it also treats
transport application through the barycentric projection
\begin{equation}
T_\epsilon(x_i)
=
\frac{1}{a_i}\sum_{j=1}^m P^\ast_{ij}y_j,
\label{eq:flash-barycentric}
\end{equation}
with the common gradient expression
\begin{equation}
\nabla_{x_i} \OT_\epsilon
=
2a_i\bigl(x_i-T_\epsilon(x_i)\bigr).
\label{eq:flash-gradient}
\end{equation}
Thus exact online/GPU Sinkhorn is the semantics-preserving engineering lane:
it improves memory traffic and GPU utilization while keeping the same entropic
transport object.  It does not, however, change the all-pairs asymptotic cost.

\section{Nystrom Kernel Sinkhorn}

The target paper \emph{Massively scalable Sinkhorn distances via the Nystrom
method} approximates the Gaussian Gibbs kernel \citep{altschuler2019nystrom}
\begin{equation}
K_{ij}
=
k_\eta(x_i,x_j)
=
\exp(-\eta\|x_i-x_j\|^2)
\label{eq:nystrom-kernel}
\end{equation}
by a low-rank factorization
\begin{equation}
\tilde K = V A^{-1}V^\top.
\label{eq:nystrom-factorization}
\end{equation}
Given landmark points
\begin{equation}
X_r=\{\tilde x_1,\ldots,\tilde x_r\}\subset X,
\label{eq:nystrom-landmarks}
\end{equation}
the factors are
\begin{equation}
V_{ij}=k_\eta(x_i,\tilde x_j),
\qquad
A_{jj'}=k_\eta(\tilde x_j,\tilde x_{j'}).
\label{eq:nystrom-factors}
\end{equation}
If $A=LL^\top$, a matrix-vector product is computed as
\begin{equation}
\tilde K v
=
V\left(L^{-\top}\left(L^{-1}(V^\top v)\right)\right).
\label{eq:nystrom-matvec}
\end{equation}
After Cholesky factorization, each multiply costs $O(nr)$ rather than
$O(n^2)$.  Sinkhorn scaling uses
\begin{equation}
\tilde P
=
D_1\tilde K D_2
=
D_1 V A^{-1}V^\top D_2.
\label{eq:nystrom-plan}
\end{equation}
For particle transport, this factorization is directly useful:
\begin{equation}
\tilde P X
=
D_1V A^{-1}V^\top D_2X.
\label{eq:nystrom-transport}
\end{equation}

The adaptive Nystrom subroutine doubles the rank until the entrywise kernel
approximation error is below a threshold.  For Gaussian kernels, the diagnostic
used in the paper is
\begin{equation}
\|K-\tilde K\|_\infty
=
1-\min_i\tilde K_{ii}.
\label{eq:nystrom-diag-error}
\end{equation}
If points lie in a ball in $\R^d$, the paper gives a worst-case bound of the
form
\begin{equation}
r^\ast(X,\eta,\epsilon')
\le
3\left(
6+\frac{53}{d}\eta R^2
+\frac{3}{d}\log\frac{2n}{\epsilon'}
\right)^d.
\label{eq:nystrom-ball-rank}
\end{equation}
This is not reassuring in high ambient dimension.  The favorable case is
low intrinsic dimension.  If the data lie on a $k$-dimensional manifold, the
bound becomes
\begin{equation}
r^\ast(X,\eta,\epsilon')
\le
c_{\Omega,\eta}
\left(\log\frac{n}{\epsilon'}\right)^{5k/2}.
\label{eq:nystrom-manifold-rank}
\end{equation}

The paper also states stability in logarithmic kernel error.  If
\begin{equation}
\|\log K-\log\tilde K\|_\infty \le \epsilon',
\label{eq:nystrom-log-error}
\end{equation}
then scaling the approximate kernel gives a plan
$\tilde P=D_1\tilde K D_2$ with marginal residual
\begin{equation}
\|\tilde P\ones-p\|_1
+\|\tilde P^\top\ones-q\|_1
\le \epsilon',
\label{eq:nystrom-marginal-residual}
\end{equation}
and controlled objective error.  For LEDH-PFPF-OT, the essential extra
diagnostic is transported-particle error,
\begin{equation}
\|\tilde P X-P_\epsilon X\|.
\label{eq:nystrom-transport-error}
\end{equation}

Nystrom Sinkhorn can help when the post-flow Gibbs kernel has small effective
rank.  It can fail when the effective rank is large, the entropy is too small,
or high ambient dimension corresponds to high intrinsic dimension.

\section{Positive-Feature Sinkhorn}

Positive-feature Sinkhorn replaces the Gibbs kernel by an inner product of
positive features \citep{scetbon2020positivefeatures}.  Choose
\begin{equation}
\phi_\theta(x)\in\R^r_{+},
\label{eq:positive-feature-map}
\end{equation}
define
\begin{equation}
k_\theta(x,y)
=
\langle\phi_\theta(x),\phi_\theta(y)\rangle,
\label{eq:positive-feature-kernel}
\end{equation}
and set the associated cost
\begin{equation}
c_\theta(x,y)
=
-\epsilon\log k_\theta(x,y).
\label{eq:positive-feature-cost}
\end{equation}
For feature matrices
\begin{align}
\xi
&=
[\phi_\theta(x_1),\ldots,\phi_\theta(x_n)]
\in(\R_+^\ast)^{r\times n},
\label{eq:positive-xi}\\
\zeta
&=
[\phi_\theta(y_1),\ldots,\phi_\theta(y_m)]
\in(\R_+^\ast)^{r\times m},
\label{eq:positive-zeta}
\end{align}
the sample kernel is
\begin{equation}
K_\theta
=
\xi^\top\zeta.
\label{eq:positive-feature-factorization}
\end{equation}
Sinkhorn matrix-vector products therefore cost $O(r(n+m))$.  The dual
objective can be written
\begin{equation}
W_{\epsilon,c_\theta}(\mu,\nu)
=
\max_{\alpha\in\R^n,\beta\in\R^m}
a^\top\alpha+b^\top\beta
-\epsilon(\xi e^{\alpha/\epsilon})^\top
\zeta e^{\beta/\epsilon}
+\epsilon.
\label{eq:positive-feature-dual}
\end{equation}
The paper also uses integral kernel representations,
\begin{equation}
k(x,y)
=
\int_U \phi(x,u)^\top\phi(y,u)\,d\rho(u),
\label{eq:positive-feature-integral}
\end{equation}
and finite approximations
\begin{equation}
\phi_\theta(x)
=
\frac{1}{\sqrt r}
(\phi(x,u_1),\ldots,\phi(x,u_r)).
\label{eq:positive-feature-finite}
\end{equation}
For particle resampling, the factored transport application is
\begin{equation}
P X
=
\diag(u)\xi^\top\zeta\diag(v)X.
\label{eq:positive-feature-transport}
\end{equation}
The caution is semantic: a feature approximation changes the cost or kernel.
The transported particle cloud and marginal residuals must be compared against
the exact entropic reference.

\section{Direct Low-Rank Coupling OT}

Low-rank kernel methods approximate $K$.  Direct low-rank OT constrains or
parametrizes the coupling $P$ itself
\citep{forrow2019factored,scetbon2021lowrank,scetbon2022lowrank,scetbon2023unbalanced}.
The nonnegative rank of a nonnegative
matrix is
\begin{equation}
\rankp(M)
=
\min\left\{
r:\;M=\sum_{k=1}^r M_k,\;
\operatorname{rank}(M_k)=1,\;M_k\ge0
\right\}.
\label{eq:nonnegative-rank}
\end{equation}
The rank-constrained coupling set is
\begin{equation}
\Pi_{a,b}(r)
=
\{P\in\Pi(a,b):\rankp(P)\le r\}.
\label{eq:low-rank-polytope}
\end{equation}
The low-rank OT problem is
\begin{equation}
\min_{P\in\Pi_{a,b}(r)}\langle C,P\rangle.
\label{eq:low-rank-ot}
\end{equation}

The factored coupling parameterization is
\begin{equation}
P
=
Q\diag(1/g)R^\top,
\label{eq:factored-coupling}
\end{equation}
where
\begin{equation}
Q\in\Pi(a,g),
\qquad
R\in\Pi(b,g),
\qquad
g\in\Delta_r.
\label{eq:factored-coupling-constraints}
\end{equation}
Equivalently,
\begin{equation}
FC_{a,b}(r)
=
\{(Q,R,g): Q\in\R_+^{n\times r},
R\in\R_+^{m\times r},
g\in(\R_+^\ast)^r,\;
Q\in\Pi(a,g),\;R\in\Pi(b,g)\}.
\label{eq:fc-set}
\end{equation}
Then
\begin{equation}
\min_{P\in\Pi_{a,b}(r)}\langle C,P\rangle
=
\min_{(Q,R,g)\in FC_{a,b}(r)}
\left\langle C,Q\diag(1/g)R^\top\right\rangle.
\label{eq:factored-low-rank-objective}
\end{equation}
Transport application is efficient:
\begin{equation}
PX
=
Q\diag(1/g)R^\top X.
\label{eq:factored-transport}
\end{equation}

A newer latent-coupling factorization introduces an additional $r\times r$
latent plan \citep{halmos2024frlc,halmos2025hierarchical}:
\begin{equation}
LC_{a,b}(r)
=
\{(Q,R,T):
Q\in\R_+^{n\times r},
R\in\R_+^{m\times r},
T\in\R_+^{r\times r},
Q\in\Pi(a,\cdot),
R\in\Pi(b,\cdot),
T\in\Pi(g_Q,g_R)
\}.
\label{eq:latent-coupling-set}
\end{equation}
The induced coupling is
\begin{equation}
P
=
Q\diag(1/g_Q)
T
\diag(1/g_R)R^\top.
\label{eq:latent-coupling}
\end{equation}
This allows mass to route through a small latent transport plan rather than
only through diagonal latent components.

The Riemannian low-rank line uses positive factors and manifold optimization
\citep{jawanpuria2026riemannian}.  For balanced linear OT, one representative
parameterization is
\begin{equation}
\Gamma
=
U\diag(U^\top\ones)^{-1}V^\top,
\label{eq:riemannian-coupling}
\end{equation}
with objective
\begin{equation}
f(U,V)=\langle C,\Gamma\rangle.
\label{eq:riemannian-objective}
\end{equation}
It also provides a rank-sufficiency diagnostic.  Let $\Gamma_r$ be the rank-$r$
solution and define the reduced gradient or dual slack matrix
\begin{equation}
R
=
\nabla_\Gamma f(\Gamma_r)
-\alpha^\ast\ones^\top
-\ones(\beta^\ast)^\top.
\label{eq:rank-slack}
\end{equation}
Then
\begin{equation}
\delta_r=\min_{ij}R_{ij}.
\label{eq:rank-certificate}
\end{equation}
Under the theorem's assumptions, $\delta_r\ge0$ certifies global optimality for
the full coupling problem; $\delta_r<0$ identifies a direction in which rank
should be increased.  Even when the exact certificate is not achieved at a
finite iterate, it is a useful diagnostic for whether the rank is too small.

Direct low-rank coupling is a promising subquadratic route because it reduces
storage and transport application.  It is also a real algorithmic change: it
solves over low-rank couplings, not the original dense entropic feasible set.

\section{Sparse, Screened, and Stabilized OT}

Sparse multiscale OT begins with the exact discrete problem
\citep{schmitzer2016sparse}
\begin{equation}
\Pi(\mu,\nu)
=
\{\pi\in P(X\times Y):
\pi(\{x\}\times Y)=\mu(x),
\pi(X\times\{y\})=\nu(y)\},
\label{eq:sparse-polytope}
\end{equation}
\begin{equation}
\min_{\pi\in\Pi(\mu,\nu)}
C(\pi),
\qquad
C(\pi)=\sum_{(x,y)\in X\times Y}c(x,y)\pi(x,y).
\label{eq:sparse-ot}
\end{equation}
The dense problem considers the full product $X\times Y$.  Sparse multiscale
methods solve a sequence of sparse subproblems while verifying global
optimality through local shielding neighborhoods.  This is relevant if
post-flow LEDH particles have genuinely local couplings.

Screenkhorn accelerates regularized OT by screening dual variables
\citep{alaya2019screenkhorn}.  The
regularized primal is
\begin{equation}
S_\eta(\mu,\nu)
=
\min_{P\in\Pi(\mu,\nu)}
\{\langle C,P\rangle-\eta H(P)\}.
\label{eq:screenkhorn-primal}
\end{equation}
With Gibbs kernel
\begin{equation}
K=\exp(-C/\eta),
\label{eq:screenkhorn-kernel}
\end{equation}
the dual can be written
\begin{equation}
\min_{u\in\R^n,v\in\R^m}
\Psi(u,v)
=
\ones_n^\top B(u,v)\ones_m
-\langle u,\mu\rangle
-\langle v,\nu\rangle,
\label{eq:screenkhorn-dual}
\end{equation}
where
\begin{equation}
B(u,v)
=
\Delta(e^u)K\Delta(e^v).
\label{eq:screenkhorn-B}
\end{equation}
The idea is to identify negligible components and solve a smaller constrained
dual problem.  As with sparse multiscale OT, the method is useful only if the
screened plan still gives acceptable transport and marginal residuals.

Stabilized sparse scaling, epsilon scaling, truncation, and coarse-to-fine
methods should be treated as a toolbox rather than a single replacement
algorithm \citep{schmitzer2019stabilized}.  They are particularly important
when $\epsilon$ is small and standard-domain $K=\exp(-C/\epsilon)$ becomes
numerically fragile.

\section{Accelerated, Greedy, and Stochastic Sinkhorn Iterations}

This family improves iteration complexity or update scheduling but does not by
itself remove dense kernel access
\citep{altschuler2017nearlinear,dvurechensky2018computational,abid2018greedy,lin2019efficient,xu2026accelerating}.
Greenkhorn updates the most violating row or column rather than all rows and
columns.  The plan still has the scaling
form
\begin{equation}
P=\diag(u)K\diag(v),
\label{eq:greenkhorn-plan}
\end{equation}
but the update schedule for components of $u$ and $v$ changes.  Accelerated
primal-dual, accelerated mirror-descent, and related methods improve
theoretical dependence on accuracy or regularization.  If each iteration still
requires dense $Kv$ and $K^\top u$, the large $N,D$ bottleneck remains.  These
methods are secondary tools to combine with online, sparse, or low-rank
backends.

\section{Stochastic and Minibatch OT}

Stochastic OT methods are important for learning losses and cases where the
measures are observed through samples
\citep{genevay2016stochastic,genevay2018samplecomplexity}.  Particle filtering
is different: at a given time step the resampling transform acts on a finite
weighted ensemble.
Replacing the full coupling by stochastic dual estimates is not automatically
a valid resampling map.

Minibatch OT methods reduce computation by solving smaller OT problems and an
outer problem over minibatches \citep{nguyen2022bomb}.  Entropic population
BoMb-OT has the form
\begin{equation}
ED^m_d(\mu,\nu)
=
\min_{\gamma\in\Pi(\mu^{\otimes m},\nu^{\otimes m})}
\mathbb{E}_{(X^m,Y^m)\sim\gamma}
\left[
  d(P_{X^m},P_{Y^m})
\right]
+\lambda\,\KL(\gamma\mid \mu^{\otimes m}\otimes\nu^{\otimes m}).
\label{eq:bomb-ot}
\end{equation}
For finite $k$ minibatches of size $m$, an outer $k\times k$ OT problem
couples minibatches.  If the inner discrepancy is Sinkhorn, cost construction
is roughly $O(m^2k^2)$.  If the inner discrepancy is sliced Wasserstein, it
can be closer to $O(m\log(m)k^2)$.

An unbalanced minibatch variant uses local unbalanced OT
\citep{nguyen2021partialminibatch}:
\begin{equation}
UOT^\phi_\tau(\mu_n,\nu_n)
=
\min_{\pi\in\R_+^{n\times n}}
\langle C,\pi\rangle
+\tau D_\phi(\pi\ones,\mu_n)
+\tau D_\phi(\pi^\top\ones,\nu_n).
\label{eq:uot}
\end{equation}
Minibatch OT is tempting for large $N$, but it is a dangerous drop-in
replacement.  A minibatch plan can be optimal inside batches and wrong for the
full ensemble.  It should be treated as a new randomized or approximate
resampling method, not as a faithful acceleration of full-ensemble transport.

\section{Sliced, Max-Sliced, Generalized Sliced, and Subspace OT}

Sliced methods attack high dimension by projecting measures to one-dimensional
distributions
\citep{kolouri2019generalized,deshpande2019maxsliced,paty2019subspace,nguyen2025sliced}.
For one-dimensional measures, the monotone transport map is
\begin{equation}
T(x)=F_\nu^{-1}(F_\mu(x)),
\label{eq:monotone-map}
\end{equation}
and
\begin{equation}
W_p^p(\mu,\nu)
=
\int_0^1
\left|
F_\mu^{-1}(z)-F_\nu^{-1}(z)
\right|^p dz.
\label{eq:oned-wasserstein}
\end{equation}
For empirical distributions this is computed by sorting projected samples.

For $\theta\in S^{d-1}$, let $\theta_{\#}\mu$ be the pushforward of $\mu$
under $x\mapsto\langle\theta,x\rangle$.  The sliced Wasserstein distance is
\begin{equation}
\SW_p^p(\mu,\nu)
=
\int_{S^{d-1}}
W_p^p(\theta_{\#}\mu,\theta_{\#}\nu)\,d\theta.
\label{eq:sliced-wasserstein}
\end{equation}
The empirical approximation uses directions $\theta_1,\ldots,\theta_L$:
\begin{equation}
\widehat{\SW}_p^p(\mu,\nu)
=
\frac{1}{L}\sum_{\ell=1}^L
W_p^p((\theta_\ell)_{\#}\mu,(\theta_\ell)_{\#}\nu).
\label{eq:empirical-sliced}
\end{equation}
Equivalently, for distributions $P_d,P_g$,
\begin{equation}
\widetilde W_p(P_d,P_g)
=
\left[
\int_{\omega\in\Omega}
W_p^p(P_d^\omega,P_g^\omega)\,d\omega
\right]^{1/p}.
\label{eq:deshpande-sw}
\end{equation}
For finite samples $D,F$ and directions $\widehat\Omega$,
\begin{equation}
\frac{1}{|\widehat\Omega|}
\sum_{\omega\in\widehat\Omega}
W_2^2(D^\omega,F^\omega)
\label{eq:finite-sliced}
\end{equation}
with each one-dimensional term computed by sorting,
\begin{equation}
W_2^2(D^\omega,F^\omega)
=
\frac{1}{|D|}\sum_i
\left\|
D^\omega_{\pi_D^\omega(i)}
-F^\omega_{\pi_F^\omega(i)}
\right\|_2^2.
\label{eq:sliced-sorting}
\end{equation}
Max-sliced Wasserstein replaces random averaging by a best projection:
\begin{equation}
\MSW_p^p(\mu,\nu)
=
\max_{\theta\in S^{d-1}}
W_p^p(\theta_{\#}\mu,\theta_{\#}\nu).
\label{eq:max-sliced}
\end{equation}

Sliced and subspace methods are highly relevant when $D$ is the dominant
problem.  They can generate projected transport maps, but a full particle
resampling transform in $\R^D$ is not the same object as a dense coupling.
They should be studied as new resampling approximations unless a precise
construction maps the projected couplings back to a validated full-state
transport.

\section{Localization and Particle-Filter Structure}

The generic OT literature does not know the state-space model.  LEDH-PFPF-OT
does.  If the state has spatial, block, factor, or conditional-independence
structure, localization may be more principled than a generic global OT
approximation.  A localized OT resampler could partition the state as
\begin{equation}
x_i=(x_i^{(1)},\ldots,x_i^{(B)}),
\label{eq:block-state}
\end{equation}
then compute block transports
\begin{equation}
P^{(b)}\in\Pi(a^{(b)},b^{(b)})
\label{eq:block-transport}
\end{equation}
or use shared particle-level couplings with local costs
\begin{equation}
C_{ij}
=
\sum_{b=1}^B C_{ij}^{(b)}.
\label{eq:block-cost}
\end{equation}
This lane should be studied alongside low-rank methods, especially for spatial
fields, macro panels, or structured latent states.

\section{Method Taxonomy}

\small
\begin{longtable}{p{0.18\textwidth}p{0.23\textwidth}p{0.14\textwidth}p{0.18\textwidth}p{0.17\textwidth}}
\toprule
Family & Representative sources & Transport object? & What it fixes & Main risk \\
\midrule
\endhead
Exact online/GPU Sinkhorn &
\citep{feydy2019geomloss,charlier2021keops,ye2026flashsinkhorn,xiao2026fastsinkhorn}
& Yes & Dense storage and GPU IO & Still all-pairs $O(N^2D)$ work \\
Nystrom Sinkhorn & \citep{altschuler2019nystrom} & Yes, factored entropic
coupling & Kernel matvecs when effective rank is small & Rank can grow badly
with dimension or small $\epsilon$ \\
Positive-feature Sinkhorn & \citep{scetbon2020positivefeatures} & Yes if
coupling application is implemented & Linear-time positive kernel matvecs &
Feature approximation changes cost/kernel \\
Low-rank coupling OT &
\citep{forrow2019factored,scetbon2021lowrank,scetbon2022lowrank,scetbon2023unbalanced,halmos2024frlc,halmos2025hierarchical,jawanpuria2026riemannian}
& Yes,
factored coupling & Coupling storage and transport application & Changes
feasible plan class \\
Sparse/multiscale/screened &
\citep{schmitzer2016sparse,schmitzer2019stabilized,alaya2019screenkhorn}
& Yes & Active pair count when coupling is local or screenable & Locality may
fail in high dimension \\
Accelerated/greedy Sinkhorn &
\citep{altschuler2017nearlinear,dvurechensky2018computational,abid2018greedy,lin2019efficient,xu2026accelerating}
& Yes &
Iteration count/convergence & Dense kernel access remains \\
Stochastic/minibatch OT &
\citep{genevay2016stochastic,genevay2018samplecomplexity,nguyen2022bomb,nguyen2021partialminibatch}
& Sometimes
partial/local plans & Training-scale memory and compute & Not a full
deterministic resampling plan \\
Sliced/subspace OT &
\citep{paty2019subspace,kolouri2019generalized,deshpande2019maxsliced,nguyen2025sliced}
&
Projected plans/maps & High-dimensional statistics and speed & Different
transport object \\
Localization/block OT & \citep{reich2013nonparametric} and model structure & Yes
if designed & Large $D$ with state locality & Requires model-specific
validation \\
\bottomrule
\end{longtable}
\normalsize

\section{Code Availability and Reuse Risk}

Source availability is a key practical filter, but it should not be confused
with default-readiness.  In this section, a \emph{valid inspected checkout}
means that the repository is locally available under
\path{.localsource/scalable_ot_code_audit}, has a readable Git commit, and the
listed implementation files were inspected.  It does not mean that the package
was installed, built, or tested.  Because BayesFilter-owned algorithmic code
defaults to TensorFlow/TFP, PyTorch, JAX, Triton, and C++ repositories below
are treated as algorithmic references unless a later plan explicitly approves a
backend exception.

\scriptsize
\setlength{\tabcolsep}{2pt}
\begin{longtable}{@{}p{0.15\textwidth}p{0.19\textwidth}p{0.12\textwidth}p{0.21\textwidth}p{0.23\textwidth}@{}}
\toprule
Repository/method & Local audit status & Backend & Transport object exposed? &
BayesFilter decision implication \\
\midrule
\endhead
FilterFlow differentiable PF
\citep{corenflos2021differentiable} &
\path{filterflow}, commit \path{5d8300b}; valid inspected checkout &
TensorFlow/TFP &
Yes.  \path{plan.py} exposes \path{transport_from_potentials} and
\path{transport}, returning a dense transport matrix from
\path{sinkhorn.py} potentials. &
Most directly compatible source for differentiable resampling semantics, but
not a scalable fix because it materializes dense $[B,N,N]$ transport. \\

POT / Python Optimal Transport \citep{flamary2021pot} &
\path{POT}, commit \path{3e10dd6}; valid inspected checkout &
Multi-backend: NumPy, PyTorch, JAX, TensorFlow, CuPy &
Yes.  Dense Sinkhorn, Greenkhorn, Screenkhorn, Nystrom lazy plans,
low-rank Sinkhorn, factored OT, and sliced plans are present in
\path{ot/bregman}, \path{ot/lowrank.py}, \path{ot/factored.py}, and
\path{ot/sliced}. &
Best broad prototype/baseline library.  It can inform TensorFlow work, but
BayesFilter still needs local wrappers and resampling-specific validation. \\

OTT-JAX \citep{cuturi2022ott} &
\path{ott-sparse}, commit \path{e241ec3}; valid inspected sparse checkout &
JAX &
Yes.  \path{SinkhornOutput} and \path{LRSinkhornOutput} expose
\path{matrix}, \path{apply}, and marginals; point-cloud and low-rank geometries
support online/factored applications. &
Excellent API reference for exact online and low-rank transport operators, but
not a default BayesFilter backend. \\

GeomLoss / KeOps online Sinkhorn
\citep{feydy2019geomloss,charlier2021keops} &
\path{geomloss}, commit \path{00e493f}; valid inspected checkout &
PyTorch/KeOps &
Yes.  The result API exposes \path{plan}, \path{lazy_plan}, and
\path{plan_operator}; point-cloud implementations expose lazy transport-plan
densities. &
Strong exact online Sinkhorn reference for avoiding dense storage.  It still
performs all-pairs interactions and would require a TensorFlow port or
backend exception. \\

FlashSinkhorn \citep{ye2026flashsinkhorn} &
\path{flash-sinkhorn}, commit \path{d538da3}; valid inspected checkout &
PyTorch/Triton &
Yes.  The API documents \path{apply_plan_vec_flashstyle} and
\path{apply_plan_mat_flashstyle}, and solver code returns shifted potentials
for GPU plan application. &
Best source for IO-aware GPU kernel design.  High porting risk into
TensorFlow, but valuable for a future fused exact-reference implementation. \\

LinearSinkhorn / positive-feature and Nystrom scripts
\citep{scetbon2020positivefeatures,altschuler2019nystrom} &
\path{LinearSinkhorn}, commit \path{01943c7}; valid inspected checkout &
PyTorch scripts &
Partial.  \path{FastSinkhorn.py} implements \path{Lin_Sinkhorn},
\path{Nys_Sinkhorn}, and feature/Nystrom helpers, but reusable particle
transport application is not exposed as a package API. &
Useful for deriving and checking positive-feature/Nystrom prototypes; not
implementation-ready for BayesFilter. \\

Schmitzer sparse/multiscale code
\citep{schmitzer2016sparse,schmitzer2019stabilized} &
\path{MultiScaleOT} commit \path{48a47aa}; \path{optimal-transport} commit
\path{1d1aaf4}; valid inspected checkouts &
C++ with Python tooling &
Yes, through sparse coupling handlers, shortcut/shielding solvers, and
stabilized sparse Sinkhorn source under \path{src/ShortCutSolver} and
\path{src/Sinkhorn}. &
Strong sparse/multiscale reference if post-flow particles have exploitable
locality.  High integration risk and likely unsuitable as first TensorFlow
prototype. \\

Mini-batch-OT / BoMb-OT
\citep{nguyen2022bomb,nguyen2021partialminibatch} &
\path{Mini-batch-OT-sparse}, commit \path{1ef5ffc}; partial checkout after
DNS/GnuTLS/HTTP2 failures &
NumPy/POT scripts &
Partially visible.  \path{ABC/utils.py} and \path{ColorTransfer/utils.py}
compute local mini-batch plans and higher-level couplings, but the checkout is
dirty/incomplete. &
Keep as concept/literature lane until a clean archive is available.  Even with
source, stochastic minibatch semantics do not directly give a single full
deterministic resampling plan. \\
\bottomrule
\end{longtable}
\setlength{\tabcolsep}{6pt}
\normalsize

The practical ranking suggested by the source audit is therefore different
from a pure literature ranking.  For a near-term BayesFilter prototype,
FilterFlow is the best source-faithful TensorFlow reference, POT and OTT-JAX
are the most useful factored/low-rank implementation references, GeomLoss and
FlashSinkhorn are the best exact online/GPU design references, and Schmitzer's
code is a specialized sparse/multiscale reference.  Mini-batch-OT should not be
used as a determining source-code factor until the checkout is clean or the
user provides an archive.

\section{Recommended Research Program}

\paragraph{Lane A: exact reference.}
Keep the current dense/streaming TensorFlow path as the correctness reference.
Study FlashSinkhorn or KeOps-style online GPU design as an implementation
benchmark \citep{charlier2021keops,ye2026flashsinkhorn,xiao2026fastsinkhorn}.
This lane answers how far exact entropic OT scales if memory traffic and kernel
fusion are fixed.  Required diagnostics include
\begin{equation}
\|P\ones-a\|_1,
\qquad
\|P^\top\ones-b\|_1,
\qquad
\|Z_{\mathrm{candidate}}-Z_{\mathrm{dense}}\|.
\label{eq:exact-diagnostics}
\end{equation}

\paragraph{Lane B: approximate entropic factored kernels.}
Prototype fixed-rank Nystrom Sinkhorn first, then compare positive-feature
Sinkhorn \citep{altschuler2019nystrom,scetbon2020positivefeatures}.  Both
should expose a factored transport application,
\begin{equation}
Z
=
\diag(b)^{-1}
\tilde P^\top X,
\label{eq:factored-entropic-Z}
\end{equation}
without dense $P$.  Primary diagnostics are effective rank, kernel
approximation error, marginal residual, and transported-particle error.

\paragraph{Lane C: structural approximate couplings.}
Study low-rank coupling OT, localized OT, and sliced/subspace transports as new
resampling methods
\citep{forrow2019factored,scetbon2021lowrank,scetbon2022lowrank,paty2019subspace,nguyen2025sliced}.
These methods should not be judged only by dense Sinkhorn parity, because some
may intentionally bias the transport to improve high-dimensional behavior.
They require downstream filtering diagnostics.

\section{Validation Ladder Before Any Default Change}

\begin{enumerate}[label=\arabic*.]
\item Small deterministic dense parity: compare transported particles, not only
OT values.
\item Marginal checks: source and target constraints with weighted inputs.
\item Positivity and non-finite checks.
\item Representation diagnostics: rank, feature dimension, sparsity, projection
count, or block locality.
\item LEDH-PFPF-OT value parity on existing fixtures.
\item Gradient smoke tests after value parity.
\item Runtime and memory ladder over $N,D$.
\item Posterior/reference diagnostics for filtering.
\item Multi-seed uncertainty-aware comparisons before ranking stochastic
methods.
\item Default-readiness review only after downstream evidence, not from paper
benchmarks alone.
\end{enumerate}

\section{Conclusions}

The answer to whether the Nystrom Sinkhorn paper \path{1812.05189} can help is
yes, but it should not be the only method considered
\citep{altschuler2019nystrom}.  Nystrom Sinkhorn is a
strong first approximate entropic transport because it yields a factored
coupling and can apply the transport to particles.  Its success depends on the
effective rank of the post-flow Gibbs kernel.

For near-term engineering, exact online/GPU Sinkhorn is the reference lane.  It
preserves the current entropic OT semantics and directly supports barycentric
transport, but it remains all-pairs.  For true large-$N,D$ scaling, the most
promising methods are structural: low-rank coupling OT, localization/block OT,
and sliced/subspace transports.  These are not mere accelerations; they are new
resampling approximations and must be validated as such.

The safest research posture is a three-lane program: exact online Sinkhorn for
reference and moderate scaling, low-rank kernel Sinkhorn for the first factored
entropic approximation, and low-rank/localized/sliced transports as explicit
new filtering methods.  No candidate should become a BayesFilter default until
it passes transported-particle parity or a declared replacement criterion,
marginal diagnostics, numerical stability checks, and downstream filtering
validation.

\section*{References}
\addcontentsline{toc}{section}{References}
\begin{thebibliography}{99}

\bibitem[Abid and Gower(2018)]{abid2018greedy}
Abid, B.~K. and Gower, R.~M. (2018).
\newblock Greedy stochastic algorithms for entropy-regularized optimal
transport problems.
\newblock \emph{arXiv:1803.01347}.

\bibitem[Alaya et~al.(2019)]{alaya2019screenkhorn}
Alaya, M.~Z., B{\'e}rar, M., Gasso, G., and Rakotomamonjy, A. (2019).
\newblock Screening Sinkhorn algorithm for regularized optimal transport.
\newblock \emph{Advances in Neural Information Processing Systems}.

\bibitem[Altschuler et~al.(2017)]{altschuler2017nearlinear}
Altschuler, J., Weed, J., and Rigollet, P. (2017).
\newblock Near-linear time approximation algorithms for optimal transport via
Sinkhorn iteration.
\newblock \emph{Advances in Neural Information Processing Systems}.

\bibitem[Altschuler et~al.(2019)]{altschuler2019nystrom}
Altschuler, J., Bach, F., Rudi, A., and Weed, J. (2019).
\newblock Massively scalable Sinkhorn distances via the Nystrom method.
\newblock \emph{Advances in Neural Information Processing Systems};
arXiv:1812.05189.

\bibitem[Charlier et~al.(2021)]{charlier2021keops}
Charlier, B., Feydy, J., Glaun{\`e}s, J.~A., Collin, F.-D., and Durif, G.
(2021).
\newblock Kernel operations on the GPU, with autodiff, without memory
overflows.
\newblock \emph{Journal of Machine Learning Research}, 22(74):1--6.

\bibitem[Corenflos et~al.(2021)]{corenflos2021differentiable}
Corenflos, A., Thornton, J., Deligiannidis, G., and Doucet, A. (2021).
\newblock Differentiable particle filtering via entropy-regularized optimal
transport.
\newblock \emph{Proceedings of the 38th International Conference on Machine
Learning}.

\bibitem[Cuturi(2013)]{cuturi2013sinkhorn}
Cuturi, M. (2013).
\newblock Sinkhorn distances: Lightspeed computation of optimal transport.
\newblock \emph{Advances in Neural Information Processing Systems}.

\bibitem[Cuturi et~al.(2022)]{cuturi2022ott}
Cuturi, M., Meng-Papaxanthos, L., Tian, Y., Bunne, C., Davis, G., and Teboul, O.
(2022).
\newblock Optimal Transport Tools (OTT): A JAX toolbox for all things
Wasserstein.
\newblock \emph{arXiv:2201.12324}.

\bibitem[Daum and Huang(2008)]{daumhuang2008}
Daum, F. and Huang, J. (2008).
\newblock Nonlinear filters with particle flow.
\newblock \emph{Signal Processing, Sensor Fusion, and Target Recognition XVII}.

\bibitem[Deshpande et~al.(2019)]{deshpande2019maxsliced}
Deshpande, I., Hu, Y.-T., Sun, R., Pyrros, A., Siddiqui, N., Koyejo, S.,
Zhao, Z., Forsyth, D., and Schwing, A.~G. (2019).
\newblock Max-sliced Wasserstein distance and its use for GANs.
\newblock \emph{Proceedings of the IEEE/CVF Conference on Computer Vision and
Pattern Recognition}.

\bibitem[Dvurechensky et~al.(2018)]{dvurechensky2018computational}
Dvurechensky, P., Gasnikov, A., and Kroshnin, A. (2018).
\newblock Computational optimal transport: Complexity by accelerated gradient
descent is better than by Sinkhorn's algorithm.
\newblock \emph{Proceedings of the 35th International Conference on Machine
Learning}.

\bibitem[Feydy et~al.(2019)]{feydy2019geomloss}
Feydy, J., S{\'e}journ{\'e}, T., Vialard, F.-X., Amari, S.-i., Trouv{\'e}, A.,
and Peyr{\'e}, G. (2019).
\newblock Interpolating between optimal transport and MMD using Sinkhorn
divergences.
\newblock \emph{Proceedings of the 22nd International Conference on Artificial
Intelligence and Statistics}.

\bibitem[Flamary et~al.(2021)]{flamary2021pot}
Flamary, R., Courty, N., Gramfort, A., Alaya, M.~Z., Boisbunon, A., Chambon,
S., Chapel, L., Corenflos, A., Fatras, K., Fournier, N., Gautheron, L.,
Gayraud, N.~T.~H., Janati, H., Rakotomamonjy, A., Redko, I., Rolet, A., Schutz,
A., Seguy, V., Sutherland, D.~J., Tavenard, R., Tong, A., and Vayer, T. (2021).
\newblock POT: Python Optimal Transport.
\newblock \emph{Journal of Machine Learning Research}, 22(78):1--8.

\bibitem[Forrow et~al.(2019)]{forrow2019factored}
Forrow, A., H{\"u}tter, J.-C., Nitzan, M., Rigollet, P., Schiebinger, G., and
Weed, J. (2019).
\newblock Statistical optimal transport via factored couplings.
\newblock \emph{Proceedings of the 22nd International Conference on Artificial
Intelligence and Statistics}.

\bibitem[Genevay et~al.(2016)]{genevay2016stochastic}
Genevay, A., Cuturi, M., Peyr{\'e}, G., and Bach, F. (2016).
\newblock Stochastic optimization for large-scale optimal transport.
\newblock \emph{Advances in Neural Information Processing Systems}.

\bibitem[Genevay et~al.(2019)]{genevay2018samplecomplexity}
Genevay, A., Chizat, L., Bach, F., Cuturi, M., and Peyr{\'e}, G. (2019).
\newblock Sample complexity of Sinkhorn divergences.
\newblock \emph{Proceedings of the 22nd International Conference on Artificial
Intelligence and Statistics}; arXiv:1810.02733.

\bibitem[Halmos et~al.(2024)]{halmos2024frlc}
Halmos, P., Liu, X., Gold, J., and Raphael, B.~J. (2024).
\newblock Low-rank optimal transport through factor relaxation with latent
coupling.
\newblock \emph{Advances in Neural Information Processing Systems}.

\bibitem[Halmos et~al.(2025)]{halmos2025hierarchical}
Halmos, P., Gold, J., Liu, X., and Raphael, B.~J. (2025).
\newblock Hierarchical refinement: Optimal transport to infinity and beyond.
\newblock \emph{arXiv:2503.03025}.

\bibitem[Jawanpuria and Mishra(2026)]{jawanpuria2026riemannian}
Jawanpuria, P. and Mishra, B. (2026).
\newblock A Riemannian approach to low-rank optimal transport.
\newblock \emph{arXiv:2606.12120}.

\bibitem[Kolouri et~al.(2019)]{kolouri2019generalized}
Kolouri, S., Nadjahi, K., Simsekli, U., Badeau, R., and Rohde, G.~K. (2019).
\newblock Generalized sliced Wasserstein distances.
\newblock \emph{Advances in Neural Information Processing Systems}.

\bibitem[Li and Coates(2017)]{li2017particle}
Li, Y. and Coates, M. (2017).
\newblock Particle filtering with invertible particle flow.
\newblock \emph{IEEE Transactions on Signal Processing}, 65(15):4102--4116.

\bibitem[Lin et~al.(2019)]{lin2019efficient}
Lin, T., Ho, N., and Jordan, M.~I. (2019).
\newblock On efficient optimal transport: An analysis of greedy and
accelerated mirror descent algorithms.
\newblock \emph{Proceedings of the 36th International Conference on Machine
Learning}.

\bibitem[Nguyen et~al.(2022)]{nguyen2022bomb}
Nguyen, K., Ho, N., Pham, T., and Bui, H. (2022).
\newblock On transportation of mini-batches: A hierarchical approach.
\newblock \emph{Proceedings of the 39th International Conference on Machine
Learning}.

\bibitem[Nguyen et~al.(2021)]{nguyen2021partialminibatch}
Nguyen, K., Nguyen, D., Pham, T., Ho, N., and Bui, H. (2021).
\newblock Improving mini-batch optimal transport via partial transportation.
\newblock \emph{arXiv:2108.09645}.

\bibitem[Nguyen(2025)]{nguyen2025sliced}
Nguyen, K. (2025).
\newblock An introduction to sliced optimal transport.
\newblock \emph{arXiv:2508.12519}.

\bibitem[Paty and Cuturi(2019)]{paty2019subspace}
Paty, F.-P. and Cuturi, M. (2019).
\newblock Subspace robust Wasserstein distances.
\newblock \emph{Proceedings of the 36th International Conference on Machine
Learning}.

\bibitem[Peyr{\'e} and Cuturi(2019)]{peyre2019computational}
Peyr{\'e}, G. and Cuturi, M. (2019).
\newblock Computational optimal transport.
\newblock \emph{Foundations and Trends in Machine Learning}, 11(5--6):355--607.

\bibitem[Reich(2013)]{reich2013nonparametric}
Reich, S. (2013).
\newblock A nonparametric ensemble transform method for Bayesian inference.
\newblock \emph{SIAM Journal on Scientific Computing}, 35(4):A2013--A2024.

\bibitem[Scetbon and Cuturi(2020)]{scetbon2020positivefeatures}
Scetbon, M. and Cuturi, M. (2020).
\newblock Linear time Sinkhorn divergences using positive features.
\newblock \emph{Advances in Neural Information Processing Systems}.

\bibitem[Scetbon et~al.(2021)]{scetbon2021lowrank}
Scetbon, M., Cuturi, M., and Peyr{\'e}, G. (2021).
\newblock Low-rank Sinkhorn factorization.
\newblock \emph{Proceedings of the 38th International Conference on Machine
Learning}.

\bibitem[Scetbon and Cuturi(2022)]{scetbon2022lowrank}
Scetbon, M. and Cuturi, M. (2022).
\newblock Low-rank optimal transport: Approximation, statistics and debiasing.
\newblock \emph{arXiv:2205.12365}.

\bibitem[Scetbon et~al.(2023)]{scetbon2023unbalanced}
Scetbon, M., Klein, M., Palla, G., and Cuturi, M. (2023).
\newblock Unbalanced low-rank optimal transport solvers.
\newblock \emph{arXiv:2305.19727}.

\bibitem[Schmitzer(2016)]{schmitzer2016sparse}
Schmitzer, B. (2016).
\newblock A sparse multi-scale algorithm for dense optimal transport.
\newblock \emph{Journal of Mathematical Imaging and Vision}, 56:238--259.

\bibitem[Schmitzer(2019)]{schmitzer2019stabilized}
Schmitzer, B. (2019).
\newblock Stabilized sparse scaling algorithms for entropy regularized
transport problems.
\newblock \emph{SIAM Journal on Scientific Computing}, 41(3):A1443--A1481.

\bibitem[Villani(2003)]{villani2003topics}
Villani, C. (2003).
\newblock \emph{Topics in Optimal Transportation}.
\newblock American Mathematical Society.

\bibitem[Xiao(2026)]{xiao2026fastsinkhorn}
Xiao, H. (2026).
\newblock Fast log-domain Sinkhorn optimal transport with warp-level GPU
reductions.
\newblock \emph{arXiv:2605.00837}.

\bibitem[Xu and Chen(2026)]{xu2026accelerating}
Xu, Z. and Chen, L. (2026).
\newblock Accelerating Sinkhorn for entropy-regularized optimal transport.
\newblock \emph{arXiv:2605.30267}.

\bibitem[Ye et~al.(2026)]{ye2026flashsinkhorn}
Ye, F.~X.-F., Li, X., Yu, A., Chang, M.-C., Chu, L., and Wertheimer, D.
(2026).
\newblock FlashSinkhorn: IO-aware entropic optimal transport on GPU.
\newblock \emph{arXiv:2602.03067}.

\end{thebibliography}

\appendix
\section{Local Source Map}
\label{sec:source-map}

\small
\begin{longtable}{p{0.30\textwidth}p{0.62\textwidth}}
\toprule
Topic & Local source \\
\midrule
\endhead
Corpus manifest & \path{.localsource/scalable_ot_survey/MANIFEST.md} \\
Code audit manifest &
\path{.localsource/scalable_ot_code_audit/MANIFEST.md} \\
Current TensorFlow annealed transport &
\path{experiments/dpf_implementation/tf_tfp/resampling/annealed_transport_tf.py} \\
LEDH/PF-PF source-form documentation &
\path{docs/chapters/ch19b_dpf_literature_survey.tex};
\path{docs/chapters/ch19c_dpf_implementation_literature.tex} \\
Cuturi, Sinkhorn distances &
\path{.localsource/dpf_ot_audit/cuturi2013_sinkhorn_distances.txt} \\
Reich, ensemble transform particle filter &
\path{.localsource/dpf_ot_audit/reich2013_nonparametric_ensemble_transform.txt} \\
Corenflos et al., differentiable OT resampling &
\path{.localsource/dpf_ot_audit/corenflos2021_differentiable_particle_filtering_eot.txt} \\
Nystrom Sinkhorn target paper &
\path{.localsource/1812.05189-src/sections/nystrom.tex};
\path{.localsource/1812.05189-src/sections/sinkhorn.tex};
\path{.localsource/1812.05189.txt} \\
Positive-feature Sinkhorn &
\path{.localsource/scalable_ot_survey/2006.07057.txt} \\
Exact online/GPU Sinkhorn &
\path{.localsource/scalable_ot_survey/1810.08278.txt};
\path{.localsource/scalable_ot_survey/2004.11127.txt};
\path{.localsource/scalable_ot_survey/2602.03067.txt};
\path{.localsource/scalable_ot_survey/2605.00837.txt} \\
Low-rank coupling OT &
\path{.localsource/scalable_ot_survey/1806.07348.txt};
\path{.localsource/scalable_ot_survey/2103.04737.txt};
\path{.localsource/scalable_ot_survey/2205.12365.txt};
\path{.localsource/scalable_ot_survey/2305.19727.txt};
\path{.localsource/scalable_ot_survey/2411.10555.txt};
\path{.localsource/scalable_ot_survey/2503.03025.txt};
\path{.localsource/scalable_ot_survey/2603.03578.txt};
\path{.localsource/scalable_ot_survey/2606.12120.txt} \\
Sparse/screened/stabilized OT &
\path{.localsource/scalable_ot_survey/1510.05466.txt};
\path{.localsource/scalable_ot_survey/1906.08540.txt};
\path{.localsource/dpf_ot_audit/schmitzer2019_stabilized_sparse_scaling.txt} \\
Accelerated/greedy Sinkhorn &
\path{.localsource/scalable_ot_survey/1705.09634.txt};
\path{.localsource/scalable_ot_survey/1802.04367.txt};
\path{.localsource/scalable_ot_survey/1803.01347.txt};
\path{.localsource/scalable_ot_survey/1901.06482.txt};
\path{.localsource/scalable_ot_survey/2605.30267.txt} \\
Stochastic/minibatch OT &
\path{.localsource/scalable_ot_survey/1605.08527.txt};
\path{.localsource/scalable_ot_survey/2102.05912.txt};
\path{.localsource/scalable_ot_survey/2108.09645.txt} \\
Sliced/subspace OT &
\path{.localsource/scalable_ot_survey/1901.08949.txt};
\path{.localsource/scalable_ot_survey/1902.00434.txt};
\path{.localsource/scalable_ot_survey/1904.05877.txt};
\path{.localsource/scalable_ot_survey/2508.12519.txt} \\
Inspected scalable OT code repositories &
\path{.localsource/scalable_ot_code_audit/filterflow};
\path{.localsource/scalable_ot_code_audit/POT};
\path{.localsource/scalable_ot_code_audit/ott-sparse};
\path{.localsource/scalable_ot_code_audit/geomloss};
\path{.localsource/scalable_ot_code_audit/flash-sinkhorn};
\path{.localsource/scalable_ot_code_audit/LinearSinkhorn};
\path{.localsource/scalable_ot_code_audit/schmitzer-MultiScaleOT};
\path{.localsource/scalable_ot_code_audit/schmitzer-optimal-transport};
\path{.localsource/scalable_ot_code_audit/Mini-batch-OT-sparse} \\
\bottomrule
\end{longtable}
\normalsize

\section{Incomplete Download Note}

An additional attempt to download a recent arXiv source produced the partial
file \path{.localsource/scalable_ot_survey/2505.06835}.  The file has a PDF
header but fails PDF metadata inspection because the trailer/xref table is
missing.  It is therefore not treated as an inspected source in this survey.

\end{document}
